\documentclass[UTF8,aspectratio=169]{ctexbeamer}
\usetheme{Madrid}
\usecolortheme{default}
\usepackage{amsmath}
\usepackage{booktabs}
\usepackage{siunitx}
\usepackage{graphicx}

\title{Quarter-Root-Only Ablation\\of the Two-Regime Penetration Model}
\subtitle{Why dropping $k_{\sqrt{}}\sqrt{t}$ improves the fit}
\author{Mie Spray Postprocessing Project}
\date{\today}

\begin{document}

\begin{frame}
  \titlepage
\end{frame}

\begin{frame}{Talk Outline}
\begin{itemize}
  \item \textbf{Background:} the original two-regime model and what it stores.
  \item \textbf{Motivation:} the $\sqrt{t}$ branch contributes essentially nothing
        in successful fits.
  \item \textbf{Ablation:} replace by a single quarter-root branch with sigmoid onset.
  \item \textbf{Results:} success rate $+28.6\%$ absolute; flagged-set median RMSE
        drops from $8.12\,$mm to $1.08\,$mm.
\end{itemize}
\end{frame}

\section{Background}

\begin{frame}{Original Two-Regime Model}
Each plume penetration $S(t)$ is fitted by
\[
S(t) = (1-w(t))\,k_{\sqrt{}}\,\sqrt{t} + w(t)\,k_{1/4}\,t^{1/4},
\quad
w(t)=\sigma\!\left(\frac{t-t_0}{s}\right)
\]

\vspace{0.5em}
Physical interpretation:
\begin{itemize}
  \item \textbf{Early phase} ($t<t_0$): momentum-driven free jet $\Rightarrow \sqrt{t}$
        (Naber--Siebers).
  \item \textbf{Late phase} ($t>t_0$): impingement / deceleration $\Rightarrow t^{1/4}$.
  \item Sigmoid $w(t)$ blends the two regimes around transition time $t_0$.
\end{itemize}

\vspace{0.5em}
Stored constants per curve (log-space):
$\log k_{\sqrt{}},\ \log k_{1/4},\ \log t_0,\ \log s$ \quad (4 parameters).
\end{frame}

\section{Motivation}

\begin{frame}{Empirical Observation: $\sqrt{t}$ branch is negligible}
\begin{center}
\includegraphics[width=\textwidth,height=0.78\textheight,keepaspectratio]%
{figures/ksqrt_ratio_distribution.png}
\end{center}
\end{frame}

\begin{frame}{Reading the Distributions}
\textbf{Left panel:} amplitude ratio at the transition time $t_0$
\[
r \;=\; \frac{k_{\sqrt{}}\,\sqrt{t_0}}{k_{1/4}\,t_0^{1/4}}
\]
\begin{itemize}
  \item Median $r \approx 2.9\times 10^{-9}$.
  \item \textbf{100\% of accepted fits} have $r<0.5$.
  \item Bulk of mass sits 8--10 decades \emph{below} the quarter branch.
\end{itemize}

\vspace{0.4em}
\textbf{Right panel:} actual amplitudes at $t_0$
\begin{itemize}
  \item Quarter branch median: $k_{1/4}\,t_0^{1/4} \approx 69\,$mm.
  \item Sqrt branch median: $\approx 2\times 10^{-7}\,$mm (numerically zero).
\end{itemize}

\vspace{0.4em}
\textbf{Implication.} The solver consistently drives $k_{\sqrt{}}$ to a degenerate
small value: the early $\sqrt{t}$ regime is not actually present in the
truncated, impingement-dominated penetration data we collect.
\end{frame}

\begin{frame}{Why this Matters for Failed Fits}
The 4-parameter model has too many degrees of freedom relative to what the
data constrains:
\begin{itemize}
  \item On clean (well-behaved) traces, the solver still converges, but with
        $k_{\sqrt{}}$ effectively dead.
  \item On flagged / hard traces, $k_{\sqrt{}}$ becomes a destructive degree of
        freedom: it inflates during fitting and \textbf{dominates the
        extrapolation} to $5\,$ms used by the masking pipeline.
  \item Result: many physically reasonable curves are killed by the
        \texttt{penetration\_far} check, even when the observed window is
        well-fitted.
\end{itemize}

\vspace{0.4em}
\textbf{Hypothesis:} a 3-parameter quarter-only model with sigmoid onset will
match the same window quality on clean data, while extrapolating safely on
flagged data.
\end{frame}

\section{Ablation}

\begin{frame}{Quarter-Only Ablation Model (\texttt{q1})}
Replace the blend by a single regime gated by a sigmoid onset:
\[
S_{q1}(t) \;=\; \sigma\!\left(\frac{t-t_0}{s}\right)\,k_{1/4}\,t^{1/4}
\]
Stored constants:
$\log k_{1/4},\ \log t_0,\ \log s$ \quad (3 parameters).

\vspace{0.5em}
\textbf{Implementation.} Run alongside the original model in
\texttt{fit\_raw\_data.py}:
\begin{itemize}
  \item \texttt{ABLATION\_QUARTER\_ONLY = True} fits both models per series.
  \item Same Huber-loss \texttt{least\_squares} solver, fresh initial guess.
  \item Equivalent masking pipeline applied to q1 (\texttt{flag\_bad\_fit\_q1}):
        finite-param checks, $t_0,s$ bounds, far-time penetration in
        $[18,\,300]\,$mm, and grouped robust-$z$ outlier detection.
  \item Success criterion: $\lnot\,\text{flag\_bad\_fit} \;\wedge\;
        \text{RMSE} < 3\,$mm \quad (symmetric across both models).
\end{itemize}
\end{frame}

\section{Results}

\begin{frame}{Aggregate Success Rate (266 folders, 23{,}495 series)}
\begin{center}
\begin{tabular}{lcc}
\toprule
& Main (4-param) & q1 (3-param) \\
\midrule
Successful series           & 14{,}406  & 21{,}142 \\
Success rate                & $61.3\%$  & $\mathbf{90.0\%}$ \\
Absolute gain               & \multicolumn{2}{c}{$+28.6$ percentage points} \\
\bottomrule
\end{tabular}
\end{center}

\vspace{0.7em}
\textbf{Per-nozzle gains} (q1 minus main, percentage points):
\begin{itemize}
  \item DS300 (highly impinged):\quad $+37.6$
  \item Nozzle 7:\quad $+28.7$
  \item Nozzle 3 (after MIN\_SERIES\_POINTS filter):\quad $+22.4$
  \item Nozzle 5:\quad $+19.6$
  \item Nozzles 1, 2, 4, 6, 8:\quad $+10$ to $+14$
\end{itemize}
\textbf{No nozzle regressed.} The improvement is largest where impingement
dominates --- consistent with the physical motivation.
\end{frame}

\begin{frame}{RMSE: Where the Real Win Lies}
\begin{center}
\begin{tabular}{lcc}
\toprule
Median RMSE                  & Main          & q1 \\
\midrule
Clean subset                 & $1.23\,$mm    & $1.27\,$mm \\
Flagged subset               & $8.12\,$mm    & $\mathbf{1.08\,\text{mm}}$ \\
\bottomrule
\end{tabular}
\end{center}

\vspace{0.7em}
Two readings:
\begin{itemize}
  \item \textbf{On clean traces} the two models are statistically indistinguishable
        --- removing $k_{\sqrt{}}$ costs nothing.
  \item \textbf{On flagged traces} q1 reduces median residual by $\approx 7.5\times$.
        The flagged set becomes \emph{usefully fitted} rather than degenerately
        extrapolated.
\end{itemize}
\end{frame}

\begin{frame}{Visual Diagnostic (typical flagged folder)}
On flagged plots the two models are visually distinguishable:
\begin{itemize}
  \item \textbf{Main (dashed)} curves diverge to $100$--$130\,$mm by $2.5\,$ms,
        driven by an inflated $\sqrt{t}$ tail outside the observation window.
  \item \textbf{q1 (dotted)} curves track the data within the window and
        plateau under $100\,$mm at long times --- the bounded $t^{1/4}$
        behavior the masking pipeline expects.
\end{itemize}
This explains the asymmetry quantified above: same data, better extrapolation
behavior, more series surviving the physical-plausibility gate.
\end{frame}

\section{Conclusion}

\begin{frame}{Conclusion and Next Steps}
\textbf{Take-aways}
\begin{itemize}
  \item The $k_{\sqrt{}}\sqrt{t}$ branch is empirically dead in our data
        (median ratio $\sim 10^{-9}$, $100\%$ below 0.5).
  \item Removing it removes a destructive degree of freedom on hard traces.
  \item Net effect: $+28.6\,\text{pp}$ success rate, $\approx 7.5\times$ smaller
        median residual on the flagged subset, and one fewer parameter to
        store and to learn downstream.
\end{itemize}

\vspace{0.5em}
\textbf{Open questions}
\begin{itemize}
  \item Should q1 become the default model for the MLP target, with the
        4-param fit retained only as diagnostic baseline?
  \item Re-derive Stage-1 / Stage-2 features from $(\log k_{1/4},\,\log t_0,\,
        \log s)$ and check whether feature instability also drops.
\end{itemize}
\end{frame}

\end{document}
